% Author: /root/g34_eightfolds, GPT-6 Astra Pro; inherited/unexposed effort.
% Run: HODGE_DIVISION, 2026-11-21.
% Status: author whole-proof read; independent whole-proof/source audits by
% /root/imprimitive and /root pass; both exact multiplier reconstructions pass.
% This section supplies structural coverage, not the character-table scan.
\subsection{Complete cyclic-normalizer coverage for projective
\texorpdfstring{$\mathrm{PSL}_3(4)$}{PSL3(4)} representations}
\label{sec:l34-normalizer-coverage}

Write $S=\mathrm{PSL}_3(4)$ and let $\widetilde S$ be its universal
perfect central extension.  The published input used below is
Breuer--Malle--O'Brien \cite[\S3.3.1 and Lemma~3.1]{D:BMO2017}.
The complete invoked argument occupies pages~5--6 of the author version
retained with this proof.  The paper identifies
\[
 Z(\widetilde S)\simeq C_4^2\times C_3,
 \qquad \operatorname{Out}(S)\simeq D_{12},
\]
and lists all thirty-six bicyclic extensions up to isoclinism.
Here a bicyclic extension has a perfect central cover of $S$ as its
derived subgroup, cyclic center in that derived subgroup, and cyclic
outer quotient.  The distinction between isomorphism and isoclinism is
handled explicitly below.  No assertion that isoclinic groups have the
same character tables or the same element orders is used.

For a semisimple operator $T$ on an $a$-dimensional complex vector space,
put
\[
 \nu(T)=a-\max_{\lambda\in\mathbf C}
                  \dim\ker(T-\lambda I).
\]
This number is unchanged under scalar multiplication or conjugacy.

\paragraph{Lemma L34--1 (cyclic extension of an invariant representation).}
Let $P\lhd B$ be finite, with $B/P=\langle bP\rangle$ cyclic of
order $d$.  Let $\rho:P\to\mathrm{GL}(W)$ be irreducible, and suppose
that an invertible $T$ satisfies
\[
 T\rho(p)T^{-1}=\rho(bpb^{-1})\quad(p\in P).
\]
Then a scalar multiple of $T$ is the image of $b$ in an irreducible
representation of $B$ restricting to $\rho$ on $P$.

\emph{Proof.}
The operator $T^d\rho(b^d)^{-1}$ commutes with $\rho(P)$, hence equals
$cI$ by Schur's lemma.  Choose $z\in\mathbf C^*$ with $z^dc=1$.
The assignments $p\mapsto\rho(p)$ and $b\mapsto zT$ satisfy the
conjugation relations and $(zT)^d=\rho(b^d)$.
Every element of $B$ has a unique expression $pb^j$ with
$p\in P$ and $0\leq j<d$, so these relations give a representation
of $B$.  It is irreducible because its restriction to $P$ is.
This argument permits any central power relation for the chosen
extension; it does not require the extension to split.\hfill$\square$

\paragraph{Lemma L34--2 (what isoclinism preserves in this setting).}
Let $A,B$ be isoclinic finite groups.  Suppose that
$Q=A'$ and $P=B'$ are perfect, that $A/Q$ and $B/P$ are cyclic,
and that these quotients act faithfully as outer automorphisms of
$Q/Z(Q)$ and $P/Z(P)$, respectively.  Let $R$ be an irreducible
representation of $A$ whose restriction to $Q$ is irreducible.
For every generator $aQ$ of $A/Q$, there exist an irreducible
representation $R_B$ of $B$ with irreducible restriction to $P$ and
an element $b\in B$ such that $R(a)$ and $R_B(b)$ differ by a scalar,
after identifying their representation spaces.  In particular their
projective eigenspace multiplicities agree.

\emph{Proof.}
An isoclinism consists of isomorphisms
\[
 \phi:A/Z(A)\longrightarrow B/Z(B),\qquad
 \psi:Q\longrightarrow P
\]
compatible with commutators.  With the convention
$[x,y]=xyx^{-1}y^{-1}$, compatibility means
$\psi([x,y])=[x',y']$ whenever $x'Z(B)=\phi(xZ(A))$ and
$y'Z(B)=\phi(yZ(A))$.
For every $q\in Q$ one has
\[
       \phi(qZ(A))=\psi(q)Z(B).                    \tag{L34.1}
\]
Indeed this holds on commutators by compatibility and the homomorphism
property of $\phi$, and every element of $Q=A'$ is a product of
commutators.

Faithfulness of the outer quotient implies $Z(A)\subset Q$ and
$Z(B)\subset P$: a central element has trivial outer action.
Equation~(L34.1) therefore induces an isomorphism $A/Q\simeq B/P$.
Choose $b$ with $bZ(B)=\phi(aZ(A))$.  Then $bP$ generates $B/P$.
For $q\in Q$, commutator compatibility gives
\[
 \psi(aqa^{-1})
   =\psi([a,q])\psi(q)
   =[b,\psi(q)]\psi(q)=b\psi(q)b^{-1}.             \tag{L34.2}
\]
Transfer $R|_Q$ to the irreducible representation
$\rho=R|_Q\circ\psi^{-1}$ of $P$.  By (L34.2), the same operator
$R(a)$ implements conjugation by $b$ on $\rho(P)$.
Lemma~L34--1 extends $\rho$ to $B$ after rescaling that operator.
Scalar multiplication preserves all eigenspace dimensions, proving
the assertion.\hfill$\square$

\paragraph{Theorem L34--3 (complete table coverage).}
Let $W$ be a nontrivial irreducible projective representation of $S$,
and let $Q\subset\mathrm{GL}(W)$ be the faithful quasisimple image
of $\widetilde S$.  Let $T$ be a finite-order operator normalizing
$Q$.  Then, up to scalar multiplication and an identification of the
representation space, $T$ occurs as an element in an irreducible
representation of one of the following thirty-six groups whose
restriction to its perfect derived subgroup is irreducible:
\[
\begin{array}{c|c|c}
\text{perfect core}&\text{cyclic outer extensions to include}
                    &\text{number including the core}\\ \hline
S&.2_1,.2_2,.2_3,.3,.6&6\\
2.S&.2_1,.2_2,.2_3&4\\
3.S&.2_1,.2_2,.2_3,.3,.6&6\\
4_1.S&.2_1,.2_2,.2_3&4\\
4_2.S&.2_1,.2_2,.2_3&4\\
6.S&.2_1,.2_2,.2_3&4\\
12_1.S&.2_1,.2_2,.2_3&4\\
12_2.S&.2_1,.2_2,.2_3&4
\end{array}                                                     \tag{L34.3}
\]
The labels are the ATLAS labels in the cited published list.  In
particular, any scalar-invariant eigenspace inequality verified on
all those irreducible restrictions applies to every compatible
normalizing operator $T$.

\emph{Proof.}
The kernel of $\widetilde S\to Q$ is central, since the projective
representation of the simple group $S$ is nontrivial and hence
faithful.  The center of $Q$ acts by scalars, and a finite scalar
subgroup is cyclic.  Thus $Q$ is a cyclic-center quotient of
$\widetilde S$.  The kernel used here is the kernel of its actual
central character; no characteristic-subgroup property is assumed.

Let $d$ be the order of the outer automorphism of $S$ induced by
$T$.  If $d=1$, choose $q\in Q$ inducing the same inner
automorphism.  Then $q^{-1}T$ centralizes the irreducible group $Q$,
so $T$ is a scalar multiple of $q$.  The core groups in (L34.3)
already cover this case by the cited classification.

Suppose $d>1$.  There is $q\in Q$ such that conjugation by $T^d$
equals conjugation by $q$ on $Q$.  To see that equality on $S$
also gives equality on $Q$, note that an automorphism of the
perfect group $Q$ inducing the identity on $Q/Z(Q)$ has the form
$x\mapsto xz(x)$ with $z:Q\to Z(Q)$ a homomorphism, hence is
the identity.  Therefore $T^dq^{-1}=cI$ by Schur's lemma.
Rescale $T$ to $T_0$ with $T_0^d=q$, and put
\[
                       A=\langle Q,T_0\rangle.
\]
This group is finite: it is the union of the $d$ cosets
$QT_0^j$, $0\leq j<d$.  Its quotient $A/Q$ is cyclic of order
exactly $d$, because its image in $\operatorname{Out}(S)$ has
that order.  Since $Q$ is perfect and $A/Q$ is abelian,
$A'=Q$.  Thus $A$ is a bicyclic extension of precisely the type
classified in \cite[Lemma~3.1]{D:BMO2017}.

Choose its isoclinism representative $B$ from (L34.3).
Both outer quotients are faithful, and the perfect derived subgroup
of $B$ is the displayed core.  Lemma~L34--2 applies to the actual
matrix representation of $A$, which is irreducible on $Q$.
It identifies the projective spectrum of $T_0$, and therefore that
of $T$, with an element in an irreducible representation of $B$
having irreducible core restriction.  This proves coverage.
It also explains exactly why one fixed table for each listed
isoclinism class suffices: central power relations are absorbed by
the scalar in Lemma~L34--1.\hfill$\square$

\paragraph{Finite multiplier check and source precision.}
The complete structural classification in Theorem~L34--3 is a
published input, not a newly proved classification.  The retained
independent program \texttt{check\_multiplier.py} checks its finite
kernel-orbit count directly from the displayed action in the source.
Write $M_4=\langle a,b\rangle\simeq C_4^2$, $c=(ab)^{-1}$.
The displayed order-six generator acts by
$a\mapsto b^3$, $b\mapsto c^3$, while the displayed $2_2$
fixes $a$ and interchanges $b,c$.  The first fixes $C_3$ and
the second inverts it.  These maps generate $D_{12}$ on
$M_4\times C_3$.
There are twenty distinct kernels of characters of this multiplier.
Their eight outer orbits have quotient orders
\[
              1,2,3,4,4,6,12,12.
\]
The order-one and order-three quotients have outer stabilizer of
order twelve; each other orbit has stabilizer of order four.
Enumerating simultaneous orbits of a character kernel and a cyclic
subgroup of its stabilizer gives exactly thirty-six pairs, with
the profile in (L34.3).  This check imports no character-table data.

There is a minor label inconsistency in the source's subsequent
centralization prose: its displayed $2_2$ action inverts
$M_4/\langle a\rangle$ and centralizes
$M_4/\langle ab^2\rangle$, whereas the following sentences describe
these centralizations oppositely.  The same wording occurs in the
official character-library example.  Our proof uses neither of
those centralization identifications: it includes both order-four
cover types, all three involution types for each, and the complete
published Lemma~3.1 list.  The independent count treats the two
order-four orbits as unordered.  This local label issue therefore
does not leave any compatible normalizer out of (L34.3).

Finally, selecting rows by irreducibility of restriction to the
perfect core is sufficient.  Each representation constructed in
Lemma~L34--1 has that property.  Including further such rows,
whether faithful or not on the core center, can only enlarge the
scan.  The numerical eigenvalue bounds and the validity of the
retained character tables are separate checks; this section makes
no independent numerical-spectrum claim.
